\scanpage{136}
\mspage{-III- \handwritten{2} -}
\begin{flushright}number \handwritten{383}\end{flushright}
Let us consider the \uline{chambers} defined by $\underline{H}$, that is, the connected components of the union of the hyperplanes $H\in\underline{H}$ (\S\ 1, number 1) ; recall that a chamber is an open convex cone (in the same place).

PROPOSITION 1.- \uline{Let} $C$ \uline{be a chamber}, $H_1,\ldots,H_m$ \uline{its walls}. \uline{The reflections} $s_{H_1},\ldots,s_{H_m}$ \uline{generate} $G$. \uline{The group} $G$ \uline{acts transitively on the set of chambers defined by} $\underline{H}$.

Let $G'$ be the subgroup of $G$ generated by $s_{H_1},\ldots,s_{H_m}$, $x$ a point of $C$, $C'$ another chamber, and $x'$ a point of $C'$. Among all the transforms $g'(x')$ $(g'\in G')$, let $y$ be one of those nearest to $x$. If $y$ were not in $C$, there would exist a wall $H_i$ of $C$ such that $x$ and $y$ lie on opposite sides of $H_i$ ; then the reflected point $s_{H_i}(y)$ would be strictly closer to $x$ than $y$ ; we therefore have $y\in C$, which shows that $G'$ acts transitively on the set of chambers defined by $\underline{H}$, and proves the second assertion.

Now let $H\in\underline{H}$. By \S\ 1, number 1, lemma 3, $H$ is a wall of a chamber $C'$ defined by $\underline{H}$. Thus, by what we have just seen, there exist an element $g'$ of $G'$ and an index $i$ $(1\handwritten{\leq}i\handwritten{\leq}m)$ such that $H=g'(H_i)$. By (1) we therefore have $s_H=g's_{H_i}{g'}^{-1}$, whence $s_H\in G'$ ; this shows that $G=G'$. Which was to be proved.

We shall henceforth fix a chamber $C$ (defined by $\underline{H}$), denote its walls by $H_1,\ldots,H_m$, and put $s_{H_i}=s_i$. Let $e_i$ be the unit vector orthogonal to $H_i$ and lying on the same side of $H_i$ as $C$. The chamber $C$ is therefore the set of points $x$ of $E$ such that
\[
\begin{array}{l@{\qquad}l}(2)&(x\handwritten{|}e_i)\handwritten{>}0\qquad\text{for all }i=1,\ldots,m.\end{array}
\]

PROPOSITION 2.- \uline{Let} $i,j$ \uline{be two distinct indices}, \uline{and} $m_{ij}$ \uline{the order of} $s_i s_j$ \uline{in} $G$. \uline{We then have} $(e_i\handwritten{|}e_j)=-\cos(\handwritten{\pi}/m_{ij})$, \uline{so that} $(e_i\handwritten{|}e_j)\handwritten{\leq}0$. \uline{If} $(e_i\handwritten{|}e_j)\handwritten{<}0$, \uline{we have} $(e_i\handwritten{|}e_j)\handwritten{\leq}-1/2$.

Let $G^{\handwritten{\unsure{o|0|\circ}}}$ be the subgroup of $G$ generated by $s_i$ and $s_j$ ; the essential subspace of $E$
% archive sheet number: 135
